\section*{Introduction}
Disorder is often represented as scattering that shortens the mean free path and reduces conductivity in conventional transport theory. In 1958, Anderson demonstrated that this picture is not always enough. His work showed that static disorder can suppress long-range electronic transport through interference of scattered electronic waves. Anderson's paper opened an interesting avenue: for a given disorder of some strength and energy, are the electronic states extended or localized? The answer depends not only on the disorder strength and energy but also on the dimensionality of the system\cite{Anderson1958}. In one dimensional systems it was shown even for weak disorder all states are localized \cite{Mott01041961,10.1098/rspa.1963.0148}. Instead of varying the amount of disorder one can consider varying the energy of eigenstates. This led to the idea of \textit{mobility edge}, the critical energy separating localized and extended states when such separation exists.\cite{https://doi.org/10.1002/pssb.19670210135} Mobility edge marks the metal-insulator transition, but the transition behaviour is the next question. Mott argued that the transition is discontinuous in 3-D system. He proposed a minimum metallic conductivity on the metallic side and a drop to zero conductivity at the mobility edge\cite{Mott1974}. This proposition was extended to 2-D system by Licciardello and Thouless, they showed $\sigma_{min}$ is universal, that is independent of length scale\cite{PhysRevLett.35.1475}. This is questioned by the scaling theory of localization, as it predicts a continuous transition in 3D and all states are localized in 2D.This debate sharpened the question of whether a genuine metallic state could persist in two dimensions.\cite{LeeRamakrishnan1985} 

Wegner developed a
scaling description based on the theory of critical phenomena and proposed that, for spatial dimensions $d>2$, the conductivity vanishes continuously at the mobility edge according to
\begin{equation}
    \sigma(E) \propto (E-E_c)^{(d-2)\nu},
\end{equation}
where $E_c$ is the mobility-edge energy and $\nu$ is the localization-length critical exponent \cite{Wegner1976}. Schuster subsequently argued for a universal conductivity jump in two dimensions, drawing on a proposed connection between localization and a disordered planar XY model \cite{AharonyImry1977,Schuster1978}. These competing predictions made the nature of the metal--insulator transition, particularly in two dimensions, a central problem in the theory of disordered electronic systems.
An important conceptual step toward resolving this problem came from Thouless, who related the conductance of a finite disordered sample to the sensitivity of its energy levels to changes in boundary conditions \cite{Thouless1974}. This provided a way to investigate localization through the behaviour of a system as its characteristic length scale increases as discussed in the next section. Building on this idea Abrahams, Anderson, Licciardello, and Ramakrishnan (AALR) developed a one-parameter scaling theory in 1979, proposing that the scale dependence of dimensionless conductance is governed by a scaling function \cite{AALR1979}. Their analysis challenged the proposed universal conductivity jump in two dimensions and predicted that there is no true metallic behaviour in that dimension: the conductance decreases gradually at first and crosses over to exponential decay at sufficiently large length scales. In dimensions $d>2$, their theory instead predicts an unstable fixed point associated with the mobility-edge transition \cite{AALR1979,LeeRamakrishnan1985}.

The analysis in this report focuses on the zero temperature scaling theory for non-interacting electrons in a static random potential. Its aim is to reconstruct AALR's argument, fill in mathematical steps omitted in the original Letter.
We first establish the idea of dimensionless conductance and remark that that is the relevant scaling parameter. With this parameter we define an scaling function and examine the consequences of the scaling function in one, two, and three dimensions.
